% 苏布拉马尼扬·钱德拉塞卡（综述）
% license CCBYSA3
% type Wiki

本文根据 CC-BY-SA 协议转载翻译自维基百科\href{https://en.wikipedia.org/wiki/Subrahmanyan_Chandrasekhar}{相关文章}

\begin{figure}[ht]
\centering
\includegraphics[width=6cm]{./figures/4962eacafc5afbdc.png}
\caption{} \label{fig_Sbl_1}
\end{figure}
苏布拉马尼扬·钱德拉塞卡（Subrahmanyan Chandrasekhar，/ˌtʃəndrəˈʃeɪkər/，音译：昌德拉谢卡；\(^\text{[4]}\)泰米尔语：சுப்பிரமணியன் சந்திரசேகர்，罗马化：Cuppiramaṇiyaṉ Cantiracēkar；1910年10月19日－1995年8月21日）\(^\text{[5]}\)是一位印裔美国理论物理学家，在恒星结构、恒星演化以及黑洞等领域作出了卓越贡献。他在职业生涯中也曾将相当一部分时间投入于流体力学研究，特别是稳定性与湍流理论，并在这些领域取得了重要成果。钱德拉塞卡与威廉·A·福勒因在“与恒星结构与演化有关的重要物理过程的理论研究”方面的贡献，共同获得1983年诺贝尔物理学奖。他对恒星演化的数学处理为现代天体物理学提供了众多关于大质量恒星及黑洞晚期演化阶段的理论模型。\(^\text{[6][7]}\)许多重要的科学概念、机构与发明均以他命名，包括著名的钱德拉塞卡极限与钱德拉X射线天文台。\(^\text{[8]}\)

钱德拉塞卡出生于英属印度晚期（英属拉杰时期），一生研究范围极广，涉及恒星结构、白矮星、恒星动力学、随机过程、辐射输运、氢负离子的量子理论、流体力学与磁流体稳定性、湍流、椭球平衡体的平衡与稳定性、广义相对论、黑洞的数学理论以及引力波碰撞理论等多个物理领域。\(^\text{[9]}\)在剑桥大学期间，他提出了一个解释白矮星结构的理论模型，首次将电子简并物质中电子速度导致的相对论性质量变化纳入计算。通过此理论，钱德拉塞卡得出结论：白矮星的质量不能超过太阳质量的1.44倍——这一极限即著名的钱德拉塞卡极限。钱德拉塞卡还修正了扬·奥尔特等人最早提出的恒星动力学模型，将银河系中恒星围绕银心旋转时受到的引力场涨落影响纳入考量。他为这一复杂的动力学问题给出了数学解——涉及二十个偏微分方程，并由此提出了一个新的物理量：动力摩擦。这种效应一方面会减缓恒星的运动速度，另一方面又有助于维持星团的整体稳定。钱德拉塞卡进一步将分析扩展至星际介质，证明银河中的气体与尘埃云分布极为不均匀，揭示了星系结构与动力学之间更深层次的联系。
\subsection{早年生活与教育}
钱德拉塞卡于1910年10月19日出生在英属印度的拉合尔（Lahore，今属巴基斯坦），出身于一个泰米尔人家庭。\(^\text{[11]}\)他的父母分别是希塔·巴拉克里希南（Sita Balakrishnan，1891–1931）与钱德拉塞卡拉·苏布拉马尼亚·艾耶尔。\(^\text{[12]}\)钱德拉塞卡出生时，其父在英属印度西北铁路局担任副审计长，因此一家居住在拉合尔。钱德拉塞卡（常被称为Chandra）有两位姐姐——拉贾拉克希米与巴拉帕尔瓦蒂；三位弟弟——维什瓦纳坦、巴拉克里希南与拉马纳坦；以及四位妹妹——萨拉达、维迪娅、萨维特丽和森达丽。他的叔父正是印度著名物理学家、诺贝尔奖得主钱德拉塞卡拉·文卡塔·拉曼。他的母亲是一位热爱学术与文学的女性，曾将易卜生的戏剧《玩偶之家》翻译成泰米尔语，被认为是启发钱德拉塞卡早年智识兴趣的重要人物。\(^\text{[13]}\)1916年，钱德拉塞卡一家从拉合尔迁往阿拉哈巴德，并于1918年定居在马德拉斯（Madras，今金奈）。

钱德拉塞卡在12岁之前一直由家中教师教授课程。\(^\text{[13]}\)在中学阶段，他的父亲教授他数学与物理，而母亲则教他泰米尔语。1922年至1925年间，他就读于马德拉斯特里普利坎的印度中学。随后，他进入马德拉斯总统学院就读，该校隶属于马德拉斯大学。1925年至1930年间，他在此完成本科课程。1929年，他受到阿诺德·索末菲一次讲座的启发，撰写了自己的第一篇论文——《康普顿散射与新统计》。\(^\text{[14]}\)1930年6月，他以优异成绩获得物理学荣誉学士学位。同年7月，钱德拉塞卡获得印度政府奖学金，前往剑桥大学攻读研究生，并被三一学院录取——这一机会由他曾通信交流的物理学家R. H. 福勒促成。在赴英旅途中，钱德拉塞卡利用航行时间推导出白矮星中简并电子气体的统计力学理论，并在福勒此前研究的基础上，加入了相对论修正项（详见后文“学术遗产”部分）。
\subsection{剑桥大学时期}
在进入剑桥大学的第一年，作为R. H. 福勒的研究生，钱德拉塞卡主要从事平均不透明度的计算，并将所得结果应用于改进简并恒星极限质量模型的构建。在英国皇家天文学会的会议上，他结识了著名天体物理学家E. A. 米尔恩。

1931年夏天（研究生第二年），他受马克斯·玻恩邀请前往哥廷根大学的理论物理研究所工作，研究主题包括不透明度、原子吸收系数及恒星光球层模型。在保罗·狄拉克的建议下，他在最后一年（1932年9月至1933年5月）前往哥本哈根理论物理研究所深造，在那里他结识了尼尔斯·玻尔。在哥本哈根期间，他与维克托·魏斯科普夫、莱昂·罗森菲尔德、乔治·普拉泽克及马克斯·德尔布吕克等人成为好友——他们当时都住在同一家旅馆。

因在简并恒星研究方面的成果，钱德拉塞卡获得了铜质奖章，并于1933年夏季获得剑桥大学哲学博士学位，其论文题为《旋转自引力多项式星体》。1933年10月9日，他被选为剑桥大学三一学院奖学金研究员，任期为1933年至1937年，成为继斯里尼瓦萨·拉马努金之后，第二位获得三一学院奖学金的印度人（间隔16年）。有趣的是，他原本确信自己无法获得该奖学金，因此已提前计划赴牛津大学在米尔恩门下继续研究，并甚至在当地租好了公寓。[14]

1934年，钱德拉塞卡访问了苏联列宁格勒（今圣彼得堡），结识了多位著名天体物理学家，包括维克托·安巴楚缅与尼古拉·亚历山德罗维奇·科济列夫，并与列夫·朗道相识。

在此期间，钱德拉塞卡也与英国著名天体物理学家阿瑟·爱丁顿爵士相识。爱丁顿起初对他的研究表现出兴趣，但在1935年1月的一次公开演讲中，却严厉批评了钱德拉塞卡的理论（详见下节“与爱丁顿的争论”与“钱德拉塞卡—爱丁顿之争”）。
\subsection{职业生涯与研究}
\subsubsection{早期职业生涯}
\begin{figure}[ht]
\centering
\includegraphics[width=6cm]{./figures/8f98ef5906b6192b.png}
\caption{} \label{fig_Sbl_2}
\end{figure}
1935年，钱德拉塞卡应哈佛天文台台长哈洛·沙普利之邀，前往美国担任为期三个月的理论天体物理学访问讲师。他于当年12月启程赴美。在哈佛期间，钱德拉塞卡以其卓越的学识与才华给沙普利留下了深刻印象，但他婉拒了对方提供的哈佛研究奖学金职位。此时，钱德拉塞卡结识了著名的荷兰天体物理观测学家赫拉德·柯伊伯，后者是当时白矮星领域的权威。柯伊伯不久前刚被奥托·斯特鲁维招募至耶基斯天文台，该台隶属于芝加哥大学，位于威斯康星州威廉斯湾，由该校校长罗伯特·梅纳德·哈钦斯主管。斯特鲁维早已听闻钱德拉塞卡的学术声誉，正考虑聘请他担任三位天体物理学教员之一，另两位候选人分别是柯伊伯与丹麦理论物理学家本特·斯特龙格伦。[14]在柯伊伯的推荐下，斯特鲁维于1936年3月邀请钱德拉塞卡前往耶基斯天文台，并向他正式提供职位。钱德拉塞卡对此极感兴趣，但最初仍婉拒邀请，返回英国。不久之后，校长哈钦斯在他返航途中特地发去电报敦促他接受职位。最终，钱德拉塞卡同意，并于1936年12月重返美国，出任芝加哥大学理论天体物理学助理教授。[14]此外，哈钦斯还曾在一起种族歧视事件中出面干预。当时，院长亨利·盖尔以种族理由否决了钱德拉塞卡参与斯特鲁维组织的课程教学的资格。对此，哈钦斯坚定地回应道：“无论如何，让钱德拉塞卡先生去教这门课。”[15]

钱德拉塞卡的整个职业生涯都在芝加哥大学度过。1941年，他被晋升为副教授，两年后的1943年，年仅33岁的他便升任为正教授。[14]1946年，普林斯顿大学向钱德拉塞卡提供了一个职位，以接替亨利·诺里斯·拉塞尔的空缺，并开出芝加哥大学薪资的两倍。为了挽留他，校长罗伯特·哈钦斯立即将他的薪资提高至与普林斯顿相同的水平，并成功说服他留在芝加哥。1952年，经恩里科·费米邀请，钱德拉塞卡被任命为莫顿·D·赫尔理论天体物理学杰出服务教授，并加入恩里科·费米研究所。
1953年，他与妻子拉莉莎·钱德拉塞卡正式加入美国国籍。[16]

1966年，美国国家航空航天局在芝加哥大学建立了天体物理与空间研究实验室。钱德拉塞卡在该大楼二层四个角落办公室之一办公（其余三位入驻者为约翰·A·辛普森、彼得·迈耶与尤金·N·帕克）。20世纪60年代末，高层公寓建成后，钱德拉塞卡居住于莱克肖尔大道4800号，晚年则搬至多彻斯特大楼5550号。
\subsubsection{与爱丁顿的争论}
从剑桥大学毕业后，钱德拉塞卡与阿瑟·爱丁顿保持着密切联系。1935年，他在英国皇家天文学会会议上发表了关于恒星结构方程的完整解。然而，爱丁顿在他报告之后立即安排了一场演讲，并在会上公开批评了钱德拉塞卡的理论。这一事件使钱德拉塞卡深感沮丧，也由此引发了一场持续多年的科学争论。爱丁顿拒绝接受恒星存在质量上限的观点，并提出了自己的替代理论模型。[17]

钱德拉塞卡随后寻求多位著名物理学家的支持，包括莱昂·罗森菲尔德、尼尔斯·玻尔与克里斯蒂安·穆勒等人，他们普遍认为爱丁顿的反驳缺乏科学依据。在整个1930年代，两人之间的紧张关系持续存在。爱丁顿在多个学术会议上公开批评钱德拉塞卡的研究，双方也多次在论文中互相对比与辩驳。最终，钱德拉塞卡于1939年完成了他关于白矮星理论的系统研究，并得到了学界广泛的赞誉。爱丁顿于1944年去世，而尽管两人曾有严重分歧，钱德拉塞卡仍多次公开表示，他深深敬佩爱丁顿，并一直视他为朋友。[17]
\subsubsection{第二次世界大战时期}
在第二次世界大战期间，钱德拉塞卡受聘于马里兰州阿伯丁试验场的弹道研究实验室工作。期间，他研究了多项与弹道学和冲击波物理相关的问题，并撰写了多份技术报告，包括：《平面激波的衰减》（1943年）《105毫米炮弹的最佳爆炸高度》《三重激波存在的条件》[18]《由弹丸反射波干涉所产生拍频波确定弹丸速度的方法》[19]《爆炸波的正常反射》[20][9]由于他在流体力学与气体动力学方面的专长，罗伯特·奥本海默曾邀请他加入洛斯阿拉莫斯的曼哈顿计划。然而，由于**安全许可审批延误，钱德拉塞卡最终未能参与该计划。据传，他当时曾短暂访问过加洛川电磁分离项目。
\subsubsection{系统化的研究哲学}
钱德拉塞卡曾写道，他从事科学研究的动力源于希望“以自己最大的能力参与科学各个领域的进步”。在他看来，贯穿其全部学术工作的首要动机是系统化。他在论文中如此阐述：“科学家所努力追求的，本质上是选择一个特定的领域、一个特定的方面或一个特定的细节，观察它能否在一个具有形式与连贯性的整体体系中找到恰当的位置；如果不能，他就必须继续探索更多信息，以帮助他实现这一目标。”[21]

钱德拉塞卡形成了一种独特的研究风格——他善于在多个物理与天体物理学领域中逐一精通、逐一系统化。因此，他的科研生涯可被划分为若干相对独立的阶段。他在每个阶段中都会深入研究某一主题，连续发表多篇论文，随后撰写一本系统总结该领域核心思想的专著，然后转向新的研究方向。这种系统化的研究节奏贯穿了他的一生：1929–1939年：研究恒星结构，尤其是白矮星理论；1939–1943年：转向恒星动力学与布朗运动理论；1943–1950年：研究辐射输运理论与氢负离子的量子理论；1950–1961年：集中于湍流理论、流体力学与磁流体稳定性；1960年代：研究椭球平衡体的平衡与稳定以及广义相对论；1971–1983年：深入发展黑洞的数学理论；1980年代后期：专注于引力波碰撞理论。[9]钱德拉塞卡这种“逐个征服、逐步系统化”的科研方式，使他在多个物理学分支中都留下了深远的理论遗产。
\subsubsection{与学生的合作}
钱德拉塞卡非常重视与学生的合作，并为此深感自豪。他曾指出，在他长达半个世纪（约1930–1980）的科研生涯中，合作者的平均年龄始终保持在30岁左右——这意味着他总是与年轻一代共同探索前沿科学。他严格要求学生在获得博士学位前，必须称呼他为“钱德拉塞卡教授”；而在博士毕业后，学生与其他同事一样，可直接称呼他为“钱德拉”。在20世纪40年代任职于耶基斯天文台期间，钱德拉塞卡每周末都会往返150英里（约240公里）驱车前往芝加哥大学授课。当时听他课的学生中有两位后来成为诺贝尔奖得主——李政道与杨振宁——而他们的获奖时间甚至早于钱德拉塞卡本人。著名天体物理学家卡尔·萨根曾回忆钱德拉塞卡的课堂氛围：“那些毫无准备、提出无关紧要问题的学生，往往会被他‘当场处决’；但凡是有价值的问题，钱德拉塞卡都会给予极为认真、细致的回应。”[22]
\subsubsection{其他学术与出版活动}
从1952年至1971年，钱德拉塞卡担任《天体物理学期刊》主编。[23]1957年，当尤金·帕克提交论文，首次提出“太阳风”的理论时，两位审稿人均予以否定。然而，钱德拉塞卡作为主编，经过审阅后认为帕克的工作在数学上没有错误，于是顶住压力，于1958年毅然发表了这篇论文。[24]这项决定后来被证明具有划时代意义。

1990年至1995年间，钱德拉塞卡致力于一个特别的研究项目——他尝试使用现代微积分语言与方法，重新阐释艾萨克·牛顿爵士的《自然哲学的数学原理》中的几何推理过程。这一工作最终形成著作《给普通读者的〈原理〉》，于1995年出版。此外，钱德拉塞卡还在引力波碰撞理论[25]与代数特殊扰动[26]等课题上进行了深入研究。
\subsection{个人生活}
\begin{figure}[ht]
\centering
\includegraphics[width=8cm]{./figures/eb6692de5b844125.png}
\caption{钱德拉塞卡与他的兄弟姐妹在印度马德拉斯（现金奈）合影。从左至右坐着的分别是：巴拉、萨维特丽、钱德拉、萨拉达、拉贾姆；站立者从左至右依次为：维迪娅、巴拉克里希南、维什瓦姆、拉马纳坦与森达丽。} \label{fig_Sbl_3}
\end{figure}
钱德拉塞卡是C. V. 拉曼的侄子，后者曾于1930年获得诺贝尔物理学奖。

1936年9月，钱德拉塞卡与拉莉莎·多拉伊斯瓦米结婚。两人最初相识于马德拉斯总统学院，当时同为在校学生。1953年，他正式成为美国归化公民。在人格与风格上，许多同事与学生认为钱德拉塞卡温和、积极、慷慨、谦逊、细致且乐于讨论；但也有人形容他在私人事务上较为内敛、严肃、急躁甚至固执，[22]对那些嘲讽或否定他研究成果的人，则极少宽恕。[27]钱德拉塞卡是一位素食主义者。[28]

1995年，钱德拉塞卡在芝加哥大学医院因心脏病发作去世，享年84岁。他早在1975年就曾经历过一次心脏病发作。[22]他的妻子拉莉莎·钱德拉塞卡比他多活了18年，于2013年9月2日逝世，享年102岁。[29]她是一位文学与西方古典音乐的热心研究者。[27]

有一次在讨论《薄伽梵歌》时，钱德拉塞卡说道：“在发言之前，我想先作一个个人声明，以免接下来的言论被误解——我认为自己是一个无神论者。”[30]他在其他多次演讲中也表达过相同立场。卡梅什瓦尔·C·瓦利曾引用他说：“我在任何意义上都不是一个宗教人士；事实上，我认为自己是个无神论者。”[31]在1987年10月6日于芝加哥大学接受凯文·克里斯丘纳斯采访时，钱德拉塞卡提到：“当然，（奥托·斯特鲁维）知道我是无神论者，但他从未在我面前提起过这个话题。”[32]出于对丈夫无神论信仰的尊重，钱德拉塞卡的妻子在家中从不陈列她随身带来的神像或宗教饰物。[33]
\subsection{奖项、荣誉与遗产}
\subsubsection{诺贝尔奖}
1983年，钱德拉塞卡因其在恒星结构与演化中关键物理过程的理论研究而获得诺贝尔物理学奖的一半。虽然他接受了这一荣誉，但他对诺贝尔委员会的引文感到失望——因为奖项仅仅提及了他早期的研究成果，而未能体现他一生的系统性贡献。他认为这在某种程度上贬低了自己毕生的科学成就。该奖项的另一半授予了威廉·A·福勒。
\subsubsection{其他奖项与荣誉}
\begin{itemize}
\item 英国皇家学会院士（1944年）[1]
\item 美国哲学学会会员（1945年）[34]
\item 美国艺术与科学院院士（1946年）[35]
\item 亨利·诺里斯·罗素讲座奖（1949年）[36]
\item 布鲁斯奖章（1952年）[37]
\item 英国皇家天文学会金质奖章（1953年）[38]
\item 美国国家科学院院士（1955年）[39]
\item 鲁姆福德奖，美国艺术与科学院颁发（1957年）[40]
\item 美国国家科学奖章（1966年）[41]
\item 印度国家勋章“莲花装勋章二级”（1968年）
\item 美国国家科学院德雷珀奖章（1971年）[42]
\item 玛丽安·斯摩鲁霍夫斯基奖章（1973年）
\item 英国皇家学会科普利奖章（1984年）
\item 戈登·J·莱因奖（1989年）
\item 美国成就学院金盘奖（1990年）[43]
\item 詹斯基讲座，于美国国家射电天文台发表
\end{itemize}
\subsubsection{学术遗产}
钱德拉塞卡最著名的科学成就之一是天体物理学中的“钱德拉塞卡极限”。该极限给出了白矮星的最大质量——约为1.44倍太阳质量，也即是恒星在超新星爆发后坍缩成中子星或黑洞所需超过的最小质量。这一极限最早由钱德拉塞卡于1930年在他从印度前往英国剑桥攻读研究生的首次航行途中计算得出。1979年，美国国家航空航天局将其“四大太空天文台”之一命名为钱德拉X射线天文台，以纪念他的贡献。该命名来自一次全球征名活动，共收到来自50个州与61个国家的6000份提名。该天文台于1999年7月23日由哥伦比亚号航天飞机发射升空并部署入轨。此外，多个重要科学名词与实体以他命名：钱德拉塞卡数：磁流体力学中的一个重要无量纲数；小行星1958 Chandra；喜马拉雅钱德拉望远镜。在《英国皇家学会院士传记回忆录》中，天体物理学家R. J. 泰勒*曾写道：“钱德拉塞卡是一位典型的应用数学大师——他的研究主要服务于天文学，而像他这样的人，或许再也不会出现了。”[1]

在职业生涯中，钱德拉塞卡共指导了45名博士生。[44]他去世后，妻子拉莉莎·钱德拉塞卡将他的诺贝尔奖金全数捐赠给芝加哥大学，用于设立“苏布拉马尼扬·钱德拉塞卡纪念奖学金”。该奖学金于2000年首次颁发，每年授予在物理系或天文学与天体物理学系博士项目中表现最为杰出的申请者。[45]此外，自2014年起，亚太物理学会联合会设立了“钱德拉塞卡等离子体物理奖”，以表彰在等离子体物理领域作出杰出贡献的科学家。[46]

钱德拉天体物理学院是一个面向高中生的教育项目，由麻省理工学院的科学家担任导师[47]，并由钱德拉X射线天文台资助[48]。该项目旨在培养青年学生对天体物理学的兴趣与研究能力。著名天体物理学家卡尔·萨根在其著作《魔鬼出没的世界》中称赞道：“我是在苏布拉马尼扬·钱德拉塞卡那里，第一次领悟到什么是真正的数学优雅。”2017年10月19日，为纪念钱德拉塞卡诞辰107周年及其“钱德拉塞卡极限”的发现，谷歌在28个国家同步发布了纪念他的Google Doodle。[49][50]

2010年，值钱德拉塞卡百年诞辰之际，芝加哥大学举办了题为“钱德拉塞卡百年纪念研讨会”的学术会议。与会的世界顶尖天体物理学家包括：罗杰·彭罗斯、基普·索恩、弗里曼·戴森、**贾扬特·纳利卡尔、拉希德·苏尼亚耶夫、G. 斯里尼瓦桑与**克利福德·威尔（Clifford Will）**等人。研讨会的研究报告随后于2011年结集出版，书名为：《**从流体流动到黑洞：献给钱德拉塞卡百年诞辰的致敬（Fluid Flows to Black Holes: A Tribute to S. Chandrasekhar on His Birth Centenary）**》[51][52][53]。
\subsection{出版著作}
\subsubsection{书籍}
\begin{itemize}
\item Chandrasekhar, S.**（2010）[1939/1958]，《**恒星结构研究导论（An Introduction to the Study of Stellar Structure）**》，纽约：多佛出版社（Dover）。ISBN 978-0-486-60413-8。
\item Chandrasekhar, S.**（2005）[1942]，《**恒星动力学原理（Principles of Stellar Dynamics）**》，纽约：多佛出版社。ISBN 978-0-486-44273-0。
\item Chandrasekhar, S.**（2011）[1950/1960]，《**辐射输运（Radiative Transfer）**》，纽约：多佛出版社。ISBN 978-0-486-60590-6。
\item Chandrasekhar, S.**（1975）[1960]，《**等离子体物理学（Plasma Physics）**》，芝加哥：芝加哥大学出版社（The University of Chicago Press）。ISBN 978-0-226-10084-5。
\item Chandrasekhar, S.**（1981）[1961]，《**流体力学与磁流体稳定性（Hydrodynamic and Hydromagnetic Stability）**》，纽约：多佛出版社。ISBN 978-0-486-64071-6。
\item Chandrasekhar, S.**（1987）[1969]，《**椭球平衡体（Ellipsoidal Figures of Equilibrium）**》，纽约：多佛出版社。ISBN 978-0-486-65258-0。
\item Chandrasekhar, S.**（1998）[1983]，《**黑洞的数学理论（The Mathematical Theory of Black Holes）**》，纽约：牛津大学出版社（Oxford University Press）。ISBN 978-0-19-850370-5。
\item Chandrasekhar, S.**（1983）[1983]，《**爱丁顿：他那个时代最杰出的天体物理学家（Eddington: The Most Distinguished Astrophysicist of His Time）**》，剑桥大学出版社（Cambridge University Press）。ISBN 9780521257466。
\item Chandrasekhar, S.**（1990）[1987]，《**真理与美：科学中的审美与动机（Truth and Beauty: Aesthetics and Motivations in Science）**》，芝加哥：芝加哥大学出版社。ISBN 978-0-226-10087-6。
\item Chandrasekhar, S.**（1995），《**给普通读者的〈原理〉（Newton's Principia for the Common Reader）**》，牛津：克拉伦登出版社（Clarendon Press）。ISBN 978-0-19-851744-3。
\item Spiegel, E. A.**（2011）[1954]，《**湍流理论：钱德拉塞卡1954年讲义（The Theory of Turbulence: Subrahmanyan Chandrasekhar’s 1954 Lectures）**》，荷兰：施普林格出版社（Springer）。ISBN 978-94-007-0117-5。
\end{itemize}
\subsubsection{注释}
\begin{itemize}
\item Chandrasekhar, S.**（1939）。〈**恒星系统的动力学 I–VIII（The Dynamics of Stellar Systems. I–VIII）**〉，《天体物理学期刊（The Astrophysical Journal）》，第90卷第1期，页1–154。Bibcode: 1939ApJ....90....1C。DOI: 10.1086/144094。ISSN 0004-637X。
\item Chandrasekhar, S.**（1943）。〈**物理学与天文学中的随机问题（Stochastic Problems in Physics and Astronomy）**〉，《现代物理评论（Reviews of Modern Physics）》，第15卷第1期，页1–89。Bibcode: 1943RvMP...15....1C。DOI: 10.1103/RevModPhys.15.1。ISSN 0034-6861。
\item Chandrasekhar, S.**（1993）。《**经典广义相对论（Classical General Relativity）**》，英国皇家学会（Royal Society）。
\item Chandrasekhar, S.**（1979）。〈**广义相对论的作用：回顾与展望（The Role of General Relativity: Retrospect and Prospect）**〉，国际天文学联合会会议论文集（Proc. IAU Meeting）。[54]
\item Chandrasekhar, S.**（1943）。《**恒星动力学新方法（New Methods in Stellar Dynamics）**》，纽约科学院（New York Academy of Sciences）。
\item Chandrasekhar, S.**（1954）。〈**瑞利散射下阳光天空的照度与偏振（The Illumination and Polarization of the Sunlit Sky on Rayleigh Scattering）**〉，《美国哲学学会汇刊（Transactions of the American Philosophical Society）》，第44卷第6期，页643–728。DOI: 10.2307/1005777。JSTOR 1005777。
\item Chandrasekhar, S.**（1983）。〈**恒星：它们的演化与稳定性（On Stars, Their Evolution and Their Stability）**，诺贝尔奖演讲（Nobel Lecture）〉，《现代物理评论（Reviews of Modern Physics）》，第56卷第2期，页137–147。斯德哥尔摩：诺贝尔基金会。DOI: 10.1103/RevModPhys.56.137。
\item Chandrasekhar, S.**（1981）。《**人类知识的新视野（New Horizons of Human Knowledge）**：在联合国教科文组织的系列公开讲座》，巴黎：联合国教科文组织出版社（Unesco Press）。
\item Chandrasekhar, S.**（1975）。〈**莎士比亚、牛顿与贝多芬：或，创造力的模式（Shakespeare, Newton, and Beethoven: Or, Patterns of Creativity）**〉，《当代科学（Current Science）》，第70卷第9期，页810–822。芝加哥大学出版。JSTOR 24099932。
\item Chandrasekhar, S.**（1973年7月）。〈**P.A.M. 狄拉克七十寿辰纪念（P.A.M. Dirac on His Seventieth Birthday）**〉，《当代物理学（Contemporary Physics）》，第14卷第4期，页389–394。Bibcode: 1973ConPh..14..389C。DOI: 10.1080/00107517308210761。ISSN 0010-7514。
\item Chandrasekhar, S.**（1947）。收录于 **Robert B. Heywood（编）**，《**心灵的工作：科学家篇（The Works of the Mind: The Scientist）**》，芝加哥大学出版社（University of Chicago Press），页159–179。OCLC 752682744。
\item Chandrasekhar, S.**（1995）。《**关于拉马努金雕像的回忆与发现（Reminiscences and Discoveries on Ramanujan’s Bust）**》，英国皇家学会出版。ASIN B001B12NJ8。
\item Chandrasekhar, S.**（1990）。《**如何探索广义相对论的物理内容（How One May Explore the Physical Content of the General Theory of Relativity）**》，美国数学学会（American Mathematical Society）。ASIN B001B10QTM。
\end{itemize}
\subsubsection{期刊论文}
钱德拉塞卡一生共发表了约 **380 篇论文**[55][1]。
他在 **1928 年**，仍为本科生时，发表了第一篇论文，主题是**康普顿效应（Compton Effect）**[56]；而他最后一篇论文发表于 **1995 年**，距他去世仅两个月，被接受发表的主题是**恒星的非径向振荡（Non-radial Oscillation of Stars）**[57]。芝加哥大学出版社（University of Chicago Press）** 将他的重要论文编为七卷本的《钱德拉塞卡精选论文集（Selected Papers of S. Chandrasekhar）》出版。
\begin{itemize}
\item Chandrasekhar, S.**（1989）。《**精选论文，第1卷：恒星结构与恒星大气（Selected Papers, Vol. 1: Stellar Structure and Stellar Atmospheres）**》，芝加哥：芝加哥大学出版社。ISBN 9780226100890。
\item Chandrasekhar, S.**（1989）。《**精选论文，第2卷：辐射输运与氢负离子（Selected Papers, Vol. 2: Radiative Transfer and Negative Ion of Hydrogen）**》，芝加哥：芝加哥大学出版社。ISBN 9780226100920。
\item Chandrasekhar, S.**（1989）。《**精选论文，第3卷：物理与天文学中的随机、统计与磁流体问题（Selected Papers, Vol. 3: Stochastic, Statistical and Hydromagnetic Problems in Physics and Astronomy）**》，芝加哥：芝加哥大学出版社。ISBN 9780226100944。
\item Chandrasekhar, S.**（1989）。《**精选论文，第4卷：等离子体物理学、流体与磁流体稳定性及张量-维里定理的应用（Selected Papers, Vol. 4: Plasma Physics, Hydrodynamic and Hydromagnetic Stability, and Applications of the Tensor-Virial Theorem）**》，芝加哥：芝加哥大学出版社。ISBN 9780226100975。
\item Chandrasekhar, S.**（1990）。《**精选论文，第5卷：相对论天体物理学（Selected Papers, Vol. 5: Relativistic Astrophysics）**》，芝加哥：芝加哥大学出版社。ISBN 9780226100982。
\item Chandrasekhar, S.**（1991）。《**精选论文，第6卷：黑洞与平面引力波碰撞的数学理论（Selected Papers, Vol. 6: The Mathematical Theory of Black Holes and of Colliding Plane Waves）**》，芝加哥：芝加哥大学出版社。ISBN 9780226101019。
\item Chandrasekhar, S.**（1997）。《**精选论文，第7卷：广义相对论中的恒星非径向振荡及其他著作（Selected Papers, Vol. 7: The Non-radial Oscillations of Stars in General Relativity and Other Writings）**》，芝加哥：芝加哥大学出版社。ISBN 9780226101040。
\end{itemize}
\subsubsection{关于钱德拉塞卡的书籍与论文}
\begin{itemize}
\item Miller, Arthur I.（2005）。《群星帝国：追寻黑洞之旅中的友情、执念与背叛》，波士顿：霍顿·米夫林出版社。ISBN 978-0-618-34151-1。
\item Srinivasan, G.（编）（1997）。《从白矮星到黑洞：钱德拉塞卡的遗产》，芝加哥：芝加哥大学出版社（。ISBN 978-0-226-76996-7。
\item Penrose, Roger（1996）。〈钱德拉塞卡、黑洞与奇点〉，《天体物理与天文学杂志（Journal of Astrophysics and Astronomy）》，第17卷第3–4期，页213–231。Bibcode: 1996JApA...17..213P。DOI: 10.1007/BF02702305。ISSN 0250-6335。
\item Parker, E.（1996）。〈钱德拉塞卡与磁流体力学〉，《天体物理与天文学杂志》，第17卷第3–4期，页147–166。DOI: 10.1007/BF02702301。ISSN 0250-6335。
\item Wali, Kameshwar C.（1991）。《钱德拉：S. 钱德拉塞卡传》，芝加哥：芝加哥大学出版社。ISBN 978-0-226-87054-0。
\item Wali, Kameshwar C.（编）（1997）。《钱德拉塞卡：传奇背后的男人——怀念钱德拉》，伦敦：帝国理工学院出版社。ISBN 978-1-86094-038-5。
\item Wali, Kameshwar C.（编）（2001）。《视角的追寻》，新加坡：世界科学出版社。ISBN 978-1-86094-201-3。
\item Wali, Kameshwar C.（编）（2020）。《S. 钱德拉塞卡：精选通信与对话》，世界科学出版社。ISBN 978-9811208324。
\item Wignesan, T.（编）（2004）。〈让群星黯然失色的人〉，《亚洲学者的亚洲（The Asianists’ Asia）》。ISSN 1298-0358。
\item Venkataraman, G.（1992）。《钱德拉塞卡与他的极限》，印度海得拉巴：大学出版社。ISBN 978-81-7371-035-3。
\item Saikia, D. J. 等（编）（2011）。《流体流向黑洞：钱德拉塞卡百年诞辰纪念论文集》，新加坡：世界科学出版社。ISBN 978-981-4299-57-2。
\item Ramnath, Radhika（编）（2012）。《S. 钱德拉塞卡：科学之人》，哈珀柯林斯出版社。ASIN B00C3EWIME。
\item Alic, Kameshwar C.；Wali, Kameshwar C.（编）（2011）。《科学自传：S. 钱德拉塞卡》，世界科学出版社。DOI: 10.1142/7686。ISBN 978-981-4299-57-2。
\item Salwi, Dilip（编）（2004）。《S. 钱德拉塞卡：学者型科学家》，Rupa出版社。ISBN 978-8129104915。
\item Pandey, Rakesh Kumar（编）（2017）。《钱德拉塞卡极限：白矮星的大小》，Lap Lambert Academic Publishing。ISBN 978-3330317666。
\end{itemize}
\subsection{参考文献}
\begin{enumerate}
\item Tayler, R. J.（1996）。〈苏布拉马尼扬·钱德拉塞卡：1910年10月19日—1995年8月21日〉，《英国皇家学会院士传记回忆录（Biographical Memoirs of Fellows of the Royal Society）》，第42卷，页80–94。DOI: 10.1098/rsbm.1996.0006。ISSN 0080-4606。S2CID 58736242。
\item “Subrahmanyan Chandrasekhar – 数学谱系计划（The Mathematics Genealogy Project）”**。网址：[www.genealogy.math.ndsu.nodak.edu](https://www.genealogy.math.ndsu.nodak.edu)。原始网页于2024年6月4日存档。
\item [1](https://www.refworld.org/legal/legislation/natlegbod/1951/en/97165)**
\item “伟大的印度人：苏布拉马尼扬·钱德拉塞卡教授”。NDTV，2014年1月26日发布。
\item Osterbrock, Donald E.（1998年12月）。〈苏布拉马尼扬·钱德拉塞卡（1910年10月19日－1995年8月21日）〉，《美国哲学学会会刊》，第142卷第4期，页658–665。美国哲学学会出版。ISSN 0003-049X。JSTOR 3152289。（需注册或订阅访问）
\item Vishveshwara, C. V.（2000年4月25日）。〈未写日记的片段：S. 钱德拉塞卡的回忆与反思〉，《当代科学》，第78卷第8期，页1025–1033。（PDF）
\item Horgan, J.（1994）。〈人物：苏布拉马尼扬·钱德拉塞卡——直面极限〉，《科学美国人》，第270卷第3期，页32–33。DOI: 10.1038/scientificamerican0394-32。ISSN 0036-8733。
\item Sreenivasan, K. R.（2019）。〈钱德拉塞卡的流体力学〉，《流体力学年度评论》，第51卷第1期，页1–24。Bibcode: 2019AnRFM..51....1S。DOI: 10.1146/annurev-fluid-010518-040537。ISSN 0066-4189。
\item O’Connor, J. J.；Robertson, E. F.〈苏布拉马尼扬·钱德拉塞卡〉，《传记》，苏格兰圣安德鲁斯大学数学与统计学院。检索日期：2012年5月21日。
\item “苏布拉马尼扬·钱德拉塞卡。NASA 星之子网站。检索日期：2017年10月19日。
\item “谁是 S. 钱德拉塞卡？”，《印度快报》，2017年10月19日。检索日期：2019年1月13日。
\item “Subramanyan Chandrasekhar 传记”，诺贝尔奖官方网站 (NobelPrize.org)。检索日期：2019年9月24日。
\item “S. 钱德拉塞卡：为何谷歌向他致敬”，Al Jazeera。检索日期：2017年10月18日。
\item Trehan, Surindar Kumar（1995）。〈苏布拉马尼扬·钱德拉塞卡 1910–1995〉（PDF），《印度国家科学院院士传记回忆录》，第23卷，页101–119。
\item 《传记回忆录》，1997年出版。DOI: 10.17226/5859。ISBN 978-0-309-05788-2。
\item “印度伟大的天体物理学家 S. 钱德拉塞卡：为什么谷歌涂鸦庆祝这位诺贝尔奖得主”，《金融快报》，2017年10月19日。检索日期：2017年10月19日。
\item Wali, Kameshwar C.（1982年10月1日）。〈钱德拉塞卡 vs. 埃丁顿——一场意料之外的对峙〉，《今日物理》，第35卷第10期，页33–40。Bibcode: 1982PhT....35j..33W。DOI: 10.1063/1.2914790。ISSN 0031-9228。
\item *Chandrasekhar, S.（1943）。《关于三重激波存在条件的研究》（报告），弹道研究实验室，阿伯丁试验场。报告编号：367。
\item Chandrasekhar, S.（1943）。《通过干涉修正频率波反射波产生的拍波来确定弹丸速度》（报告），弹道研究实验室，阿伯丁试验场。报告编号：365。
\item Chandrasekhar, S.（1943）。《105毫米炮弹的最佳爆炸高度》（PDF 报告），陆军弹道研究实验室，阿伯丁试验场。报告编号：BRL-MR-139。原始PDF文档于2019年11月29日存档。
\item 《心灵的工作》，第176页，罗伯特·B·海伍德（Robert B. Heywood）编，芝加哥大学出版社（University of Chicago Press），1947年。
\item Wali, Kameshwar C.（1991）。《钱德拉：S. 钱德拉塞卡传》，芝加哥：芝加哥大学出版社，第9页。ISBN 978-026870540。OCLC 21297960。
\item Helmut A. Abt（1995年12月1日）。〈讣告——苏布拉马尼扬·钱德拉塞卡〉，《天体物理学期刊》（，第454卷，页551。Bibcode: 1995ApJ...454..551A。DOI: 10.1086/176507。ISSN 0004-637X。
\item Parker, E. N.（1958年11月1日）。〈行星际气体与磁场的动力学〉，《天体物理学期刊》，第128卷，页664。Bibcode: 1958ApJ...128..664P。DOI: 10.1086/146579。ISSN 0004-637X。
\item Chandrasekhar, S.; Xanthopoulos, B. C.（1985）。〈与流体运动耦合的引力波的一些精确解〉，《伦敦皇家学会学报A：数学与物理科学》，第402卷（1823期），页205–224。Bibcode: 1985RSPSA.402..205C。DOI: 10.1098/rspa.1985.0115。S2CID 120942390。
\item Chandrasekhar, S.（1984）。〈关于黑洞的代数特殊扰动（〉，《伦敦皇家学会学报A：数学与物理科学》，第392卷（1802期），页1–13。Bibcode: 1984RSPSA.392....1C。DOI: 10.1098/rspa.1984.0021。S2CID 122585164。
\item Wali, Kameshwar C.（编）（1997）。《S. 钱德拉塞卡：传奇背后的人*》，伦敦：帝国理工学院出版社，第107页。ISBN 978-1860940385。OCLC 38847561。
\item Sullivan, Walter（1995年8月22日）。〈苏布拉马尼扬·钱德拉塞卡逝世，享年84岁；诺贝尔奖得主揭示了“白矮星”〉，《纽约时报》。检索日期：2020年9月13日。
\item “诺贝尔奖得主之妻拉莉莎·钱德拉塞卡逝世，享年102岁”，《印度教徒报》，2013年9月7日。
\item Wali, Kameshwar C.（编）（1997）。《S. 钱德拉塞卡：传奇背后的人》，伦敦：帝国理工学院出版社，1997年1月1日版。ISBN 978-1860940385。
\item Wali, Kameshwar C.（1991）。《钱德拉：钱德拉塞卡传》，芝加哥大学出版社，第304页。ISBN 9780226870557。引述自钱德拉塞卡本人：“我在任何意义上都不算宗教人士；事实上，我认为自己是一个无神论者。”
\item “与 S. 钱德拉塞卡博士的访谈”，美国物理学会。原文于2015年2月2日存档。检索日期：2010年1月13日。
\item Jacobsen, Knut A.（2023）。《印度教侨民》，牛津大学出版社（Oxford University Press），第184页。ISBN 978-0-19-886769-2。
\item “美国哲学学会会员记录”。网址：[search.amphilsoc.org](https://search.amphilsoc.org)。
\item “苏布拉马尼扬·钱德拉塞卡”，美国艺术与科学院，2023年2月9日。
\item “资助、奖项与荣誉”，美国天文学会。原文于2010年1月24日存档。检索日期：2011年2月24日。
\item “凯瑟琳·沃尔夫·布鲁斯金质奖章历届获奖者”，太平洋天文学会（Astronomical Society of the Pacific）。原文于2011年7月21日存档。检索日期：2011年2月24日。
\item “英国皇家天文学会金质奖章得主名单”（，英国皇家天文学会。原文于2011年5月25日存档。检索日期：2011年2月24日。
\item “S. 钱德拉塞卡”，美国国家科学院官方网站（[www.nasonline.org）。](http://www.nasonline.org）。)
\item “伦福德奖历届获奖者”，美国艺术与科学院。原文于2012年9月27日存档。检索日期：2011年2月24日。
\item “总统国家科学奖：获奖者详情”，美国国家科学基金会，网址：[www.nsf.gov](https://www.nsf.gov)。
\item “亨利·德雷珀奖章”，美国国家科学院（。原文于2013年1月26日存档。检索日期：2011年2月24日。
\item “美国成就学院金盘奖得主”，美国成就学院，网址：[www.achievement.org](https://www.achievement.org)。
\item Singh, Virendra（2011年10月26日）。〈S. 钱德拉塞卡：他的生命与科学〉，《共鸣》），第16卷第10期，页960。DOI: 10.1007/s12045-011-0094-0。ISSN 0971-8044。S2CID 119945333。
\item “学术项目 | 芝加哥大学天文学与天体物理系”，网址：[astrophysics.uchicago.edu](https://astrophysics.uchicago.edu)。
\item “奖项信息”，亚太物理学会等离子体物理分会（Association of Asia Pacific Physical Societies Division of Plasma Physics），网址：[aappsdpp.org](https://aappsdpp.org)。
\item Hartman, Mark; Ashton, Peter; Porro, Irene; Ahmed, Shakib; Kol, Simba.〈钱德拉天体物理学院〉，麻省理工学院开放课程网（MIT OpenCourseWare）。检索日期：2017年10月20日。
\item “钱德拉天体物理学院 – Chandra博客 – 最新钱德拉新闻”，网址：[chandra.harvard.edu](https://chandra.harvard.edu)。
\item “S. 钱德拉塞卡107岁诞辰纪念（S. Chandrasekhar's 107th Birthday）”，经由 [www.google.com](https://www.google.com)。
\item Rajamanickam Antonimuthu**（2017年10月18日）。〈S. 钱德拉塞卡谷歌涂鸦〉，原文于2021年12月11日存档，经由 YouTube。
\item “KPTC 活动视频 – 学术演讲”，芝加哥大学网站（。检索日期：2019年1月13日。
\item “纪念苏布拉马尼扬·钱德拉塞卡诞辰100周年：钱德拉塞卡百年研讨会2010 – 芝加哥”，VideoLectures.NET。检索日期：2019年1月13日。
\item “国家科学基金会奖项查询：奖项编号#1039863 – 钱德拉塞卡百年研讨会；美国伊利诺伊州芝加哥；2010年10月16–17日”，美国国家科学基金会（NSF），网址：[www.nsf.gov](https://www.nsf.gov)。检索日期：2019年1月13日。
\item Chandrasekhar, S.（1980）。〈广义相对论在天文学中的作用——回顾与展望〉，《天文学要论》，第5卷，页45–61。Bibcode: 1980HiA.....5...45C。DOI: 10.1017/S1539299600003749。ISSN 1539-2996。
\item Chandrasekhar, S.（1996）。〈S. 钱德拉塞卡的出版作品〉（PDF），《天体物理学与天文学杂志》，第17卷，印度科学院，页269。Bibcode: 1996JApA...17..269C。检索日期：2017年5月15日。
\item 〈“与恒星内部有关的康普顿效应的热力学研究”〉（PDF），《印度物理学杂志》，第3卷，页241–250。hdl:10821/487。
\item Chandrasekhar, Subrahmanyan; Ferrari, Valeria（1995年8月8日）。〈恒星的非径向振荡：V. 牛顿型恒星的完全相对论处理〉，《皇家学会学报A：数学、物理与工程科学》，第450卷（1939期），英国皇家学会出版，页463–475。Bibcode: 1995RSPSA.450..463C。DOI: 10.1098/rspa.1995.0094。ISSN 1364-5021。S2CID 120769457。
\end{enumerate}
\subsection{延伸阅读}
\begin{itemize}
\item Struve, Otto**（1952年4月1日）。〈授予S.钱德拉塞卡博士布鲁斯金质奖章〉，《太平洋天文学会刊》，第64卷（377期），页55。Bibcode: 1952PASP...64...55S。DOI: 10.1086/126422。ISSN 0004-6280。S2CID 119668926。
\item Parker, E. N.（1997）。〈苏布拉马尼扬·钱德拉塞卡，1910–1995〉（PDF），《美国国家科学院院士传记回忆录》，第72卷，页28–49。
\end{itemize}
\textbf{讣告}
\begin{itemize}
\item Devorkin, David H.（1996年1月1日）。〈讣告：苏布拉马尼扬·钱德拉塞卡，1910–1995〉，《美国天文学会公报》，第28卷（4期），页1448。Bibcode: 1996BAAS...28.1448D。
\item McCrea, W.（1996年4月1日）。〈讣告：苏布拉马尼扬·钱德拉塞卡〉，《天文台》，第116卷，页121–124。Bibcode: 1996Obs...116..121M。ISSN 0029-7704。
\item 评论：Cronin, J. W.（1998年2月1日）。〈苏布拉马尼扬·钱德拉塞卡〉，《天文台》，第118卷，页24。Bibcode: 1998Obs...118...24C。ISSN 0029-7704。
\item Garstang, R. H.（1997年2月1日）。〈苏布拉马尼扬·钱德拉塞卡（1910–1995）〉，《太平洋天文学会刊》，第109卷，页73–77。Bibcode: 1997PASP..109...73G。DOI: 10.1086/133864。ISSN 0004-6280。S2CID 123095503。
\item Spruit, H. C.（1996年3月1日）。〈引文索引上天文学家的“增长曲线”〉，《英国皇家天文学会季刊》，第37卷，页1–9。Bibcode: 1996QJRAS..37....1S。ISSN 0035-8738。
\end{itemize}
\subsection{外部链接}
\begin{itemize}
\item 《伟大的印度人：苏布拉马尼扬·钱德拉塞卡教授 —— 钱德拉塞卡在芝加哥的最后一次采访视频。
\item 音频 – 苏布拉马尼扬·钱德拉塞卡（1988）：〈广义相对论的创立及其卓越性。
\item 音频 – Cain / Gay（2010）：《天文学播客：钱德拉塞卡》。
\item 美国国家科学院传记。
\item 哈佛大学关于钱德拉塞卡的专题网页。
\item 苏布拉马尼扬·钱德拉塞卡。
\item 苏布拉马尼安·钱德拉谢卡（另一拼写形式）。
\item 布鲁斯金质奖章页面。
\item 口述历史访谈稿：1977年5月17日，与苏布拉马尼扬·钱德拉塞卡的第一场访谈，美国物理学会，尼尔斯·玻尔图书馆与档案馆。
\item 口述历史访谈稿：1977年5月18日，第二场访谈，美国物理学会，尼尔斯·玻尔图书馆与档案馆。
\item 口述历史访谈稿：1977年10月31日，第三场访谈，美国物理学会，尼尔斯·玻尔图书馆与档案馆。
\item 口述历史访谈稿：1987年10月6日，钱德拉塞卡晚年的访谈，美国物理学会，尼尔斯·玻尔图书馆与档案馆。
\item 数学家族谱项目。
\item 康考迪亚大学1988年6月授予荣誉学位致辞，存档于康考迪亚大学档案与记录管理处。
\item 《辐射传输》免费PDF版，可于Archive.org获取。
\item 《苏布拉马尼扬·钱德拉塞卡论文指南（1913–2011）》，存于芝加哥大学特藏研究中心。
\item 诺贝尔奖官方网站上的苏布拉马尼扬·钱德拉塞卡页面。
\item 可在维基数据上编辑相关条目。
\end{itemize}